\section[LiDAR konfigurazioa]{Sentsorearen konfigurazioa eta ezkerreko pareta}\label{sec:geometria}
URDF ereduaren arabera, LiDAR sentsorea robotaren aurrealdearekiko $180^\circ$ biratuta dago ($\text{yaw} = \pi$), hau da, atzerantz orientatuta. Horregatik, nodoak laserraren irakurketak berrantolatzen ditu (\texttt{scan[i - 1500]}), datuetako norabideak robotaren norabideekin bat etor daitezen. Zuzenketa horren ondoren, $0^\circ$ robotaren aurrealdea da, eta $+90^\circ$ ezkerraldea.

Hasierako posiziotik pistako geziei jarraituz abiatzean, kanpoko pareta ezkerrean geratzen zaigu. Paretaren angelua kalkulatzeko, beraz, aurrealdearen eta ezkerraldearen arteko laser-irakurketak erabiltzen ditugu, hau da, $0^\circ$ eta $+90^\circ$ artekoak. Horietatik $0{,}5\text{ m}$ baino distantzia txikiagoa dutenak soilik hartzen ditugu ($d_i < D_{\max} = 0{,}5\text{ m}$), robotaren ondoan dagoen paretaren puntuekin lan egiteko.

\section[Angeluaren kalkulua]{Paretaren angeluaren kalkulua eta zeinuen portaera}\label{sec:erregresioa}
\noindent\textit{Enuntziatuko galdera: Azaldu nola kalkulatzen duzun robota eta paretaren arteko angelua.}\par\vspace{1pt}

LiDAR sentsoreak izpi bakoitzean bi datu ematen dizkigu: izpiaren norabidea, $\beta_i$, eta norabide horretan dagoen oztoporainoko distantzia, $d_i$. Datu horiek zuzenean ez dute paretaren malda erakusten. Horregatik, lehenengo urratsean irakurketa bakoitza robotaren inguruko puntu bihurtzen dugu:
\vspace{-0.4em}
\[
  x_i = d_i \cos\beta_i, \qquad y_i = d_i \sin\beta_i
\]
\vspace{-0.4em}
$x_i$ koordenatuak puntua robotaren aurrean zenbat dagoen adierazten du, eta $y_i$ koordenatuak ezkerrerantz dagoen distantzia. Puntu guztiak plano berean jarri ondoren, \texttt{numpy.polyfit} erabiliz ahalik eta ondoen doitzen zaien zuzena kalkulatzen dugu: $y = c_1 x + c_0$. Zuzen horrek paretaren norabidea laburbiltzen du. $c_1$ maldaren arkutangentea hartuta, robotaren eta paretaren arteko $\theta$ angelua lortzen dugu. Robota paralelo joatea nahi dugunez, nahi den angelua $0^\circ$ da; ondorioz, kontrol-errorea kalkulatutako angelua bera da:
\vspace{-0.4em}
\[
  \theta = \arctan(c_1), \qquad e = \theta - 0{,}0 = \theta
\]
\vspace{-0.4em}
Interpretazioa zuzena da: $\theta$ zero inguruan badago, robota paretarekiko paralelo doa. $\theta$ negatiboa bada ($c_1 < 0$), robota ezkerreko paretarantz orientatuta dago ($y$ koordenatu txikiagoak aurrerago) eta eskuinera biratu behar du ($\omega < 0$); $\theta$ positiboa bada ($c_1 > 0$), paretatik urruntzen ari da eta ezkerrera biratu behar du ($\omega > 0$). Horrela, laser-puntuetatik lortutako angeluak zuzenean erabakitzen du biraketaren noranzkoa.

\section[Kontrol-legeak]{Kontrol-legeak eta abiadura lineala}\label{sec:kontrol-legeak}
PID kontrolagailuak errorearen hiru alderdi erabiltzen ditu. P osagaiak une horretako erroreari erantzuten dio. Horrez gain, I osagaiak errorea denboran pilatu den aztertzen du; D osagaiak, berriz, errorea zein abiaduratan aldatzen ari den hartzen du kontuan. Hiru ekarpenak batuta lortzen da robotari bidaltzen zaion biraketa-abiadura.

P kontrol hutsean, biraketa-abiadura angelu-errorearekiko proportzionala da: $\omega = K_p \cdot \theta$. $K_p$ txikia bada, robotak poliki zuzentzen du norabidea; $K_p$ handia bada, zuzenketa bortitzagoa da. Abiadura linealerako bi aukera alderatu ditugu: abiadura konstantea eta biraketaren arabera moteltzen den abiadura:
\vspace{-0.4em}
\[
  v = \frac{v_0}{1 + K_V \cdot |\omega|}, \qquad v_0 = 0{,}2\text{ m/s}
\]
\vspace{-0.4em}
$K_V = 0$ denean, robota $0{,}2\text{ m/s}$-ko abiadura konstantean doa. $K_V = 2$ denean, $|\omega|$ handitu ahala $v$ txikitzen da; beraz, robotak kurbetan motelago egiten du aurrera. \texttt{kp}, \texttt{ki} eta \texttt{kd} ROS parametro gisa jasotzen dira. $v_0$, $K_V$, $D_{\max} = 0{,}5\text{ m}$ eta $I_{\max} = 0{,}3\text{ rad}\cdot\text{s}$, aldiz, kodean definitutako konstanteak dira.

\section[PID kodean]{PID osagaiak kodean nola gehitu ditugun}\label{sec:pid-inplementazioa}
\noindent\textit{Enuntziatuko galdera: Nola gehitu duzue kodean osagai Integral eta Deribatiboa:}\par\vspace{1pt}

P kontrolari I eta D osagaiak \texttt{follow\_wall()} metodoan gehitu dizkiogu. Ziklo bakoitzean, aurreko neurketatik igarotako denbora kalkulatzen dugu. Ondoren, errorearen metaketa eguneratzen dugu I osagaia lortzeko, eta uneko errorea aurrekoarekin alderatzen dugu D osagaia kalkulatzeko. Azken biraketa-agindua hiru ekarpenen batura da:
\vspace{-0.6em}
\begin{verbatim}
dt = (self.timestamps[-1] - self.timestamps[-2]
      if len(self.timestamps) > 1 else 0.0)
self.error_sum = float(np.clip(self.error_sum + error * dt, -I_MAX, I_MAX))
derivative = (error - self.prev_error) / dt if dt > 0.0 else 0.0
self.prev_error = error
w = self.kp * error + self.ki * self.error_sum + self.kd * derivative
v = V0 / (1.0 + KV * abs(w))
\end{verbatim}
\vspace{-0.6em}
$\Delta t$ simulazio-denborako bi neurketen arteko diferentzia da (\texttt{/clock}-etik jasotakoa). Horrela, integralak eta deribatuak ziklo bakoitzean benetan igarotako denbora erabiltzen dute. Gainera, errore metatua $[-0{,}3, +0{,}3]\text{ rad}\cdot\text{s}$ tartera mugatzen dugu (\texttt{np.clip}), kurbetan gehiegi pilatu ez dadin. Deribatua $\Delta t > 0$ denean soilik kalkulatzen da, zeroz zatitzea saihesteko.
